Yes: for all sufficiently large $n$, every graph $G$ with $\omega(G),\alpha(G)\le10\log n$ contains an induced cycle on at least $c\log n$ vertices, for an absolute $c>0$. A cycle of length at least four is both nontrivial and regular of degree two. Thus the slightly stronger cycle conclusion suffices. Logarithm bases only change the constants; below all logarithms have base two.

Pr\"omel and R\"odl proved the stronger logarithmic universality theorem in \emph{Non-Ramsey graphs are $c\log n$-universal}, JCTA 88 (1999), 379--384, \url{https://doi.org/10.1006/jcta.1999.2972}. We reconstruct its consequence quantitatively from the following two explicit inputs, as stated and used in Fox--Sudakov, \emph{Induced Ramsey-type theorems}, Section 1.1 and Theorem 1.1, \url{https://kam.mff.cuni.cz/~matousek/cla/fox-sudakov-induced-ramsey.pdf}. These structural theorems are external inputs, not independently established here.
